\documentclass[10pt,a4paper]{article}

\usepackage[margin=22mm]{geometry}
\usepackage{amsmath,amssymb}
\usepackage{graphicx}
\usepackage{booktabs}
\usepackage{caption}
\usepackage{subcaption}
\usepackage{xcolor}
\usepackage{titlesec}
\usepackage{enumitem}
\usepackage{fancyhdr}
\usepackage{siunitx}
\usepackage[hidelinks]{hyperref}

\definecolor{accent}{HTML}{16407A}
\definecolor{lightgrey}{HTML}{F2F4F7}

\titleformat{\section}{\normalfont\Large\bfseries\color{accent}}{\thesection}{0.6em}{}
\titleformat{\subsection}{\normalfont\large\bfseries}{\thesubsection}{0.6em}{}

\pagestyle{fancy}
\fancyhf{}
\lhead{\small\color{accent}VARUNA \textbar{} Autonomous Underwater Reconnaissance and Assessment}
\rhead{\small\thepage}
\renewcommand{\headrulewidth}{0.4pt}

\captionsetup{font=small,labelfont={bf,color=accent}}
\sisetup{detect-all}

% Generated by eval/make_metrics_tex.py. Do not edit by hand.

\newcommand{\misDuration}{893}
\newcommand{\misDurationMin}{14.9}
\newcommand{\misPath}{573}
\newcommand{\misCoverage}{95.3}
\newcommand{\misNavMean}{0.15}
\newcommand{\misNavMax}{0.29}
\newcommand{\misNavFinal}{0.26}
\newcommand{\misSigma}{0.28}
\newcommand{\misDvl}{96.4}
\newcommand{\misSoundings}{12,924,685}
\newcommand{\misPings}{1391}
\newcommand{\misEnvelope}{2.93}
\newcommand{\siteCurrent}{2.4}
\newcommand{\siteSSC}{3.2}
\newcommand{\misWall}{287}
\newcommand{\misDriftPct}{0.04}
\newcommand{\scourPierTwoDepth}{3.01}
\newcommand{\scourPierTwoDepthTruth}{3.40}
\newcommand{\scourPierTwoVol}{289}
\newcommand{\scourPierTwoVolTruth}{321}
\newcommand{\scourPierTwoRec}{89}
\newcommand{\scourPierThreeDepth}{1.76}
\newcommand{\scourPierThreeDepthTruth}{2.10}
\newcommand{\scourPierThreeVol}{108}
\newcommand{\scourPierThreeVolTruth}{124}
\newcommand{\scourPierThreeRec}{84}
\newcommand{\scourMeanAbsErr}{0.37}
\newcommand{\scourMeanRec}{86}
\newcommand{\mapRmse}{0.28}
\newcommand{\mapBias}{+0.04}
\newcommand{\detClusters}{26}
\newcommand{\detPrimary}{3}
\newcommand{\detPrimaryTotal}{4}
\newcommand{\detMeanErr}{2.15}
\newcommand{\detMaxErr}{3.56}
\newcommand{\detDebrisHit}{1}
\newcommand{\detDebrisTot}{9}
\newcommand{\detVehicleLargeHit}{2}
\newcommand{\detVehicleLargeTot}{2}
\newcommand{\detVehicleSmallHit}{1}
\newcommand{\detVehicleSmallTot}{1}
\newcommand{\detBoulderHit}{3}
\newcommand{\detBoulderTot}{14}
\newcommand{\detStructureHit}{0}
\newcommand{\detStructureTot}{1}
\newcommand{\yoloValMapFifty}{99.3}
\newcommand{\yoloValMap}{78.9}
\newcommand{\yoloValP}{98.3}
\newcommand{\yoloValR}{98.9}
\newcommand{\yoloTestMapFifty}{99.0}
\newcommand{\yoloTestMap}{78.1}
\newcommand{\yoloTestP}{97.4}
\newcommand{\yoloTestR}{98.2}
\newcommand{\yoloClasses}{11}
\newcommand{\simRealRecall}{0.0}
\newcommand{\simRealScenes}{\textit{n/a}}
\newcommand{\simRealInstances}{\textit{n/a}}
\newcommand{\simRealConfReal}{\textit{n/a}}
\newcommand{\synthMapFifty}{0.4}
\newcommand{\synthMap}{0.1}
\newcommand{\synthP}{74.2}
\newcommand{\synthR}{2.2}
\newcommand{\synthTrain}{2200}
\newcommand{\synthSelfMapFifty}{84.6}
\newcommand{\synthSelfMap}{61.1}
\newcommand{\hilNavMean}{0.107}
\newcommand{\hilNavMax}{0.212}
\newcommand{\hilTicks}{17,884}
\newcommand{\hilSimTime}{894}
\newcommand{\hilCrc}{0}
\newcommand{\boardComputeMean}{34.6}
\newcommand{\boardComputeNn}{36.8}
\newcommand{\boardWaitMean}{7.4}
\newcommand{\boardBudget}{50}
\newcommand{\boardMargin}{1.4}
\newcommand{\boardLoopNn}{49.3}
\newcommand{\espSent}{17,884}
\newcommand{\espEchoes}{17,879}
\newcommand{\espMatched}{17,878}
\newcommand{\espMismatch}{1}
\newcommand{\espLag}{2.3}
\newcommand{\espMs}{3.3}
\newcommand{\espMsNn}{3.8}
\newcommand{\cadLength}{1074}
\newcommand{\cadDia}{180}
\newcommand{\cadSlender}{6.0}
\newcommand{\cadMidbody}{606}
\newcommand{\cadWetted}{0.554}
\newcommand{\cadSkin}{4}
\newcommand{\cadVolume}{0.0282}
\newcommand{\cadHullVolume}{0.0229}
\newcommand{\cadMass}{28.0}
\newcommand{\cadNetBuoy}{1.96}
\newcommand{\cadBG}{27.6}
\newcommand{\cadBallast}{6.3}
\newcommand{\cadParts}{52}
\newcommand{\cadVolErr}{0.001}
\newcommand{\cadLengthM}{1.07}
\newcommand{\cadDiaM}{0.18}
\newcommand{\cadBGAssumed}{85}
\newcommand{\structDepth}{50}
\newcommand{\structPressure}{0.490}
\newcommand{\structSkin}{4}
\newcommand{\structFrameSpacing}{200}
\newcommand{\structHoop}{10.8}
\newcommand{\structStrength}{250}
\newcommand{\structSfYield}{23.2}
\newcommand{\structCollapseUn}{66}
\newcommand{\structCollapseSt}{482}
\newcommand{\structSfBuckle}{9.6}
\newcommand{\structFrameGain}{7.3}
\newcommand{\structPylonArm}{214}
\newcommand{\structPylonStress}{10.3}
\newcommand{\structPylonSf}{24}
\newcommand{\hydroCf}{0.00399}
\newcommand{\hydroForm}{1.133}
\newcommand{\hydroDragHull}{1.3}
\newcommand{\hydroDragApp}{37.6}
\newcommand{\hydroDragTotal}{38.9}
\newcommand{\hydroDragSim}{32}
\newcommand{\hydroDragRatio}{1.21}
\newcommand{\hydroAppShare}{97}
\newcommand{\hydroAddAxial}{6.1}
\newcommand{\hydroAddLateral}{25.9}
\newcommand{\hydroEnvSim}{2.93}
\newcommand{\hydroEnvCad}{2.69}
\newcommand{\hydroMargin}{1.12}
\newcommand{\hydroSlender}{6.05}
\newcommand{\bomTotal}{32,825}
\newcommand{\bomLakh}{27.6}
\newcommand{\bomAcoustics}{23,710}
\newcommand{\bomAcousticsPct}{72}
\newcommand{\bomRest}{9115}
\newcommand{\bomRate}{84}
\newcommand{\enBattery}{1500}
\newcommand{\enHotel}{55}
\newcommand{\enMissionWh}{295}
\newcommand{\enMissionPct}{25}
\newcommand{\enMissions}{4.1}
\newcommand{\enPropMean}{1132}
\newcommand{\enPropPeak}{1993}
\newcommand{\enTotalMean}{1187}
\newcommand{\enHoursSite}{1.2}
\newcommand{\enHoursCalm}{17.8}
\newcommand{\enEta}{0.60}
\newcommand{\enRatio}{15}
\newcommand{\redSurgeIntact}{339}
\newcommand{\redEnvIntact}{2.69}
\newcommand{\redNeed}{277}
\newcommand{\redSingleOk}{4}
\newcommand{\redPairsOk}{6}
\newcommand{\redPairsTotal}{28}
\newcommand{\redWorstPct}{50}
\newcommand{\redDegraded}{1.86}
\newcommand{\redThrustNeeded}{196}
\newcommand{\redWorstSurge}{170}
\newcommand{\redWorstEnv}{1.83}
\newcommand{\redHeaveIntact}{480}
\newcommand{\detContacts}{24.4}
\newcommand{\detFalseAlarms}{15.4}
\newcommand{\detFalseAlarmsSd}{1.5}
\newcommand{\detFalsePerHa}{10.0}
\newcommand{\detSurveyHa}{1.54}
\newcommand{\detAllFound}{3}
\newcommand{\detRuns}{5}
\newcommand{\xdRealSelfFifty}{99.0}
\newcommand{\xdRealSelf}{78.1}
\newcommand{\xdSimSelfFifty}{84.6}
\newcommand{\xdSimSelf}{61.1}
\newcommand{\xdRealToSimFifty}{0.8}
\newcommand{\xdRealToSim}{0.1}
\newcommand{\xdSimToRealFifty}{0.4}
\newcommand{\xdSimToReal}{0.1}
\newcommand{\repRuns}{5}
\newcommand{\repCompleted}{5}
\newcommand{\repAllTargets}{3}
\newcommand{\repCoverMean}{95.3}
\newcommand{\repCoverStd}{0.2}
\newcommand{\repNavMeanMean}{0.322}
\newcommand{\repNavMeanStd}{0.161}
\newcommand{\repNavMaxMean}{0.663}
\newcommand{\repNavMaxStd}{0.293}
\newcommand{\repDvlMean}{96.0}
\newcommand{\repDvlStd}{0.3}
\newcommand{\repMapRmseMean}{0.285}
\newcommand{\repMapRmseStd}{0.001}
\newcommand{\repDetErrMean}{1.80}
\newcommand{\repDetErrStd}{0.40}
\newcommand{\repScourRecMean}{86.1}
\newcommand{\repScourRecStd}{0.1}
\newcommand{\deNFourZeroCoco}{56.4}
\newcommand{\deNFourZeroSim}{45.8}
\newcommand{\deNEightZeroCoco}{76.8}
\newcommand{\deNEightZeroSim}{64.1}
\newcommand{\deNOneSixZeroCoco}{89.3}
\newcommand{\deNOneSixZeroSim}{83.1}
\newcommand{\deNThreeTwoZeroCoco}{96.4}
\newcommand{\deNThreeTwoZeroSim}{91.8}
\newcommand{\deNSixFourZeroCoco}{97.5}
\newcommand{\deNSixFourZeroSim}{97.9}
\newcommand{\deDiffNFourZero}{-10.5}
\newcommand{\deDiffNEightZero}{-12.7}
\newcommand{\deDiffNOneSixZero}{-6.2}
\newcommand{\deDiffNThreeTwoZero}{-4.6}
\newcommand{\deDiffNSixFourZero}{+0.4}
\newcommand{\deBestGain}{+0.4}
\newcommand{\deBestN}{640}


% Figures are included only if present, so the document still builds
% while an experiment that produces one is still running.
\newcommand{\keyfig}[3]{%
  \IfFileExists{#1}{%
    \begin{figure}[htbp]\centering
    \includegraphics[width=\linewidth]{#1}
    \caption{#2}\label{#3}
    \end{figure}}{}}

\begin{document}

\begin{titlepage}
\centering
\vspace*{1.4cm}
{\color{accent}\rule{\linewidth}{1.2pt}}\\[4mm]
{\Huge\bfseries VARUNA}\\[3mm]
{\Large Autonomous Underwater Reconnaissance and Assessment\\ for Post-Flood Disaster Response}\\[3mm]
{\color{accent}\rule{\linewidth}{1.2pt}}\\[10mm]

{\large Technical Design and Verification Report}\\[3mm]
{\normalsize Simulation, autonomy and acoustic perception for zero-visibility,\\ high-current Indian river conditions}\\[14mm]

\begin{minipage}{0.86\linewidth}
\small
\textbf{Summary.} Post-flood search and bridge foundation inspection in India
depend on divers who cannot safely work in monsoon flood conditions: the water
is opaque with suspended sediment, the current exceeds human capability, and
submerged structural debris makes entry unsafe. This report presents a
complete, verified simulation and autonomy stack for an underwater vehicle
designed for exactly those conditions, together with quantitative results from
an end to end autonomous mission over a modelled post-collapse river site.

Three results are established. First, the commercial platform normally assumed
for this role cannot perform it: an open frame vehicle of that class holds
station only to \SI{0.96}{\metre\per\second} against a design current of
\SI{\siteCurrent}{\metre\per\second}. Second, a purpose designed faired
vehicle raises that envelope to \SI{\misEnvelope}{\metre\per\second} and
completes a \SI{\misPath}{\metre} autonomous survey with a mean navigation
error of \SI{\misNavMean}{\metre} without any external positioning. Third,
submerged vehicles and structural debris are located from the resulting
bathymetric map with no training data at all, at a mean localisation error of
\SI{\detMeanErr}{\metre}.
\end{minipage}

\vfill
{\small Prepared for RakshaTech Synapse 2026\\
Technology Innovation Hub of IIT Delhi (IHFC), Department of Science and Technology\\
Broad technology category 1: Autonomous and Robotic Platforms}
\end{titlepage}

\tableofcontents
\newpage

% ======================================================================
\section{Problem and operational context}

\subsection{The gap}

India experiences recurrent monsoon flooding that destroys road and rail
crossings and sweeps vehicles into rivers. Two capabilities are needed
immediately afterwards, and neither is currently available at acceptable risk:

\begin{enumerate}[leftmargin=*,itemsep=2pt]
\item \textbf{Search.} Locating submerged vehicles and casualties in turbid,
fast moving water.
\item \textbf{Structural assessment.} Measuring scour around bridge
foundations to decide whether remaining spans are safe.
\end{enumerate}

Both are performed today by divers. In flood conditions this is close to
impossible. The reference scenario used throughout this report is the collapse
of the Savitri river bridge on NH-66 at Mahad, Maharashtra, on 2 August 2016,
when a British era masonry arch bridge failed during monsoon flood and carried
two state transport buses and several private vehicles into the river.
Recovery took days, for precisely these reasons.

Optical imaging is not a candidate. At the suspended sediment concentrations
of a monsoon flood, modelled here at \SI{\siteSSC}{\gram\per\litre}, visibility
is effectively zero. The system is therefore built around acoustic perception
throughout.

\subsection{Statutory driver}

Beyond emergency response, IRC:SP:35 requires periodic underwater inspection
of bridge foundations, and the Dam Safety Act 2021 mandates inspection of
submerged structures. Scour depth is the quantity that determines whether a
pier is at risk. Measuring it currently requires divers or reservoir drawdown.
In the system described here it falls out of a routine survey automatically,
which converts an occasional, hazardous, qualitative inspection into a
repeatable quantitative measurement.

\keyfig{../results/figures/f0_architecture.png}
{System architecture. The vehicle acts only on estimated quantities; ground
truth is used solely to score the result afterwards.}{fig:arch}

% ======================================================================
\section{Why the standard platform does not work}

Proposals in this space normally assume an open frame commercial ROV of the
BlueROV2 class. The first engineering result of this work is that such a
vehicle cannot perform the mission, and the reason is not control tuning but
installed thrust against hydrodynamic drag.

For steady motion the surge drag is $X_{u}v + X_{u|u|}v|v|$. Setting this
equal to the largest pure surge wrench the thruster layout can produce before
any single thruster saturates gives the maximum current the vehicle can hold
station against. For the open frame baseline this is
\SI{0.96}{\metre\per\second}. The modelled site runs at
\SI{\siteCurrent}{\metre\per\second} at the surface. The vehicle is swept away,
and no controller can prevent it.

\keyfig{../results/figures/f5_vehicle.png}
{Left: surge drag against current speed for both vehicles, with the saturated
thrust limit of each. The intersection is the maximum holdable current. The
open frame baseline is out of envelope at the design condition by a factor of
more than two. Right: thruster utilisation over the reference mission for the
purpose designed vehicle.}{fig:vehicle}

\subsection{The design response}

Two changes recover the envelope, and a third is forced by the first:

\begin{itemize}[leftmargin=*,itemsep=2pt]
\item \textbf{Faired hull.} Axial quadratic drag is reduced by roughly a
factor of four relative to an open frame. Drag is deliberately left high
broadside, which also gives useful passive weathervaning into the flow.
\item \textbf{Larger thrusters and longer moment arms.} Force and torque share
one thrust budget, so moment arms are sized to buy attitude authority without
consuming the surge thrust that holds the vehicle in the current.
\item \textbf{Fixed stabilising fins.} A faired hull is directionally
unstable. The added mass asymmetry between axial and transverse directions
produces a destabilising Munk moment proportional to
$(Z_{\dot w} - X_{\dot u})\,u_r w_r$ whenever the hull moves at an angle of
attack. Buoyancy restoring is far too weak to counter it. Fins are therefore a
requirement, not a refinement, and they only stabilise when the vehicle flies
nose first.
\end{itemize}

\begin{table}[htbp]\centering\small
\caption{Vehicle comparison at the design condition.}
\begin{tabular}{lcc}
\toprule
& \textbf{Open frame baseline} & \textbf{VARUNA-1} \\
\midrule
Mass & \SI{13.5}{\kilogram} & \SI{28.0}{\kilogram} \\
Axial quadratic drag $X_{u|u|}$ & 141 & 32 \\
Thrust per unit & \SI{51}{\newton} & \SI{120}{\newton} \\
Stabilising fins & none & yes \\
Maximum holdable current & \SI{0.96}{\metre\per\second} & \SI{\misEnvelope}{\metre\per\second} \\
Station keeping at \SI{1.96}{\metre\per\second} & swept, \SI{14.0}{\metre} error & holds, \SI{0.40}{\metre} RMSE \\
\bottomrule
\end{tabular}
\end{table}

\subsection{Consequences for the mission plan}

Because the hull is only stable flying nose first, and because a broadside
turn at the design current is unrecoverable within the available sway thrust,
two constraints propagate into the autonomy layer: every commanded heading
points into the flow, and survey legs run along the flow axis rather than
across it. This is an example of vehicle physics dictating mission structure,
and it is captured directly in the planner.

% ======================================================================
\section{Acoustic simulation}

The core of this work is a physics based forward looking sonar simulator.
It forms images from the active sonar equation applied per range and bearing
cell, rather than texturing a depth buffer, and it runs on CPU alone.

\subsection{Image formation}

For a resolution cell at slant range $r$ insonified by beam $b$,

\begin{equation}
\mathrm{EL}(r,b) = \mathrm{SL} - 2\,\mathrm{TL}(r) + \mathrm{BS}(r,b),
\qquad
\mathrm{TL}(r) = 20\log_{10} r + \alpha r ,
\end{equation}

with backscattering strength following a Lambert like angular law
$\mathrm{BS} = \mu_{\text{material}} + 10 n \log_{10}(\cos i)$ plus a specular
term for smooth hard faces. The absorption coefficient
$\alpha = \alpha_{\text{water}}(f) + \alpha_{\text{sediment}}(C)$ uses the pure
water term of the Francois and Garrison formulation, which at
\SI{750}{\kilo\hertz} gives \SI{61}{\decibel\per\kilo\metre}. The sediment term
represents excess attenuation from suspended solids in flood water and is an
engineering approximation, stated as such.

Three properties of real imagery are reproduced explicitly, because they are
what a perception model actually keys on:

\begin{enumerate}[leftmargin=*,itemsep=2pt]
\item \textbf{Acoustic shadow.} Ray casting resolves occlusion, so every object
throws a shadow whose length encodes its height. Shadow geometry is the primary
recognition cue in sonar imagery.
\item \textbf{Elevation ambiguity.} A forward looking sonar collapses the
vertical aperture: all scatterers within the vertical beamwidth at a given
slant range fold into one pixel. This is reproduced by casting many elevation
rays per beam and summing them under the vertical beam pattern.
\item \textbf{Speckle.} Coherent imaging gives Rayleigh distributed amplitude,
so intensity is Gamma distributed with shape equal to the number of looks.
\end{enumerate}

\keyfig{../results/figures/f0b_sonar_pipeline.png}
{Sonar image formation. Every stage is a physical effect rather than a
rendering convenience.}{fig:pipeline}

\subsection{Footprint treatment}

Each elevation ray represents a solid angle, not a point. Its intersection with
a surface spans a finite range interval, and its energy is deposited across
that interval. For an elevation beam of angular width $\Delta\theta$ striking a
surface at grazing angle $g$, the along range extent is
\begin{equation}
L = r\,\Delta\theta\,\cot g, \qquad \sin g = |\hat d \cdot \hat n| .
\end{equation}
Deriving $L$ from local geometry rather than from neighbouring rays matters at
object edges. There the adjacent ray may strike a surface tens of metres
further away, and a midpoint rule would smear that ray's energy across the
very shadow the object is casting. This was observed directly during
verification and corrected.

\subsection{Verification}

The model is verified against physical expectations rather than by inspection.
Ray primitive intersections are checked against analytic values to
\num{e-9}. For a steel target on a silt bed the model produces a
\SI{+12.7}{\decibel} target return and a \SI{-43.6}{\decibel} acoustic shadow
behind it, with the seabed return falling smoothly with range. One ping costs
roughly \SI{150}{\milli\second} on eight CPU cores, which is faster than the
\SI{2}{\hertz} the vehicle actually uses.

\keyfig{../results/figures/real_vs_sim.png}
{Top: real ARIS Explorer 3000 captures from the public marine debris dataset.
Bottom: frames from this simulator. Both show the characteristic bright
target return with a dark shadow behind it, speckle texture, and range
dependent seabed brightness.}{fig:realvssim}

% ======================================================================
\section{Autonomy stack}

\subsection{Navigation}

There is no GPS underwater. The pose estimate comes from a twelve state
extended Kalman filter carrying position, body velocity, attitude and gyro
bias, corrected by a Doppler velocity log, a depth cell and the inertial
attitude solution. Sensor models include the failure behaviour that actually
limits the solution, in particular loss of DVL bottom lock over soft sediment,
which the research literature identifies as a principal failure mode.

Horizontal position is never directly observed. It is dead reckoned from DVL
velocity, so its error grows without bound. The filter therefore carries and
reports its own position covariance, which the mission layer can use to decide
when a position reset against a mapped structure is required.

\subsection{Control}

A cascaded controller maps position error to a desired ground velocity, then
to a body wrench. The current is fed forward as a \emph{force}, being the
thrust needed to balance hull drag at the estimated flow speed. Adding it to
the velocity setpoint instead would be incorrect, because the inner loop
regulates ground referenced velocity from the DVL, and inflating that setpoint
commands the vehicle to fly downstream rather than hold position. This error
was made and corrected during development, and it is worth stating because it
is not obvious.

The ambient current is estimated without any flow sensor, by inverting the
drag model against the applied thrust and the measured ground velocity. Over
the reference mission this recovers the true current to within
\SI{0.01}{\metre\per\second} in steady conditions.

Thrust allocation is prioritised. When demand exceeds what the thrusters can
deliver, translational force is preserved and attitude torque is scaled back,
because losing station in fast water is unrecoverable whereas a slow attitude
response is not. Uniform scaling of the whole wrench does the opposite: during
development a large yaw demand was found to collapse delivered surge thrust
from \SI{340}{\newton} to \SI{67}{\newton}, which alone caused the vehicle to
be swept away.

% ======================================================================
\section{Mapping, scour and target detection}

Every ping yields returns across the swath, which are projected through the
estimated vehicle pose and accumulated into a gridded bathymetric map. Because
the map is built in the estimated frame, navigation drift propagates into it,
so all mapping results below measure the whole chain rather than the sonar
alone.

\subsection{Scour quantification}

The undisturbed bed level around a pier is taken as the median of an annulus
outside the scour hole. Depth is the deepest excursion below that datum and
volume is the integral of the depression. If the pier surroundings are not
sufficiently sounded the result is reported as unmeasured rather than as a
spuriously shallow value.

\keyfig{../results/figures/f4_scour.png}
{Radial bed profiles around the inspected piers, mapped against ground truth,
and the resulting scour depth recovery.}{fig:scour}

\subsection{Detection without training data}

The bathymetric residual above a smooth estimate of the bare bed isolates
objects standing proud of the riverbed. The bed is estimated by grey scale
morphological opening, which removes compact objects while following broad
bathymetry including scour hollows. The structuring element must exceed the
largest object of interest: a window shorter than an \SI{11}{\metre} bus
absorbs the bus into the bed estimate and hides it entirely.

This gives a purely geometric detector that needs no labelled sonar at all.
That matters because no public training data exists for submerged Indian
buses, and it is the target class that matters most.

\keyfig{../results/figures/f7_detection.png}
{Left: height above the fitted bare bed surface. Right: detections against
ground truth, with green lines joining each matched pair.}{fig:detect}

% ======================================================================
\section{Learned acoustic perception}

A YOLOv8s detector was trained on the public marine debris forward looking
sonar dataset, captured with an ARIS Explorer 3000 and comprising
\num{2035} real images across \yoloClasses{} classes. On the held out test
split it reaches \SI{\yoloTestMapFifty}{\percent} mAP at IoU 0.5,
\SI{\yoloTestP}{\percent} precision and \SI{\yoloTestR}{\percent} recall.

This number should be read with care. The dataset is a controlled water tank
collection with a consistent background, so it is an easier problem than field
sonar and near saturated performance is expected. It establishes that the
perception pipeline is sound; it does not by itself predict field performance.

\subsection{Domain transfer}

Two transfer directions were examined. Running the real trained detector
directly on simulated open river frames gives essentially no detections. The
diagnosis is that the gap is dominated by scene content and capture geometry,
not by acoustic physics: the tank captures are dense with walls, rig
structure and cabling, at a short range window and a near grazing look, none
of which resembles an open riverbed survey.

The operationally relevant direction is the reverse, and it is the one the
system depends on, because there is no labelled sonar for Indian flood
scenarios and there will not be one before deployment. Results for a detector
trained only on simulator output and evaluated on the real held out test split
are reported in Section~\ref{sec:results}.

% ======================================================================
\section{Results}\label{sec:results}

\subsection{Reference mission}

A single autonomous mission over the modelled Savitri site: deploy, acquire
bottom lock, lawnmower search of the box, orbit inspection of two piers,
return and report. No operator input after launch.

\keyfig{../results/figures/f1_mission_overview.png}
{Mission track over site bathymetry, coloured by mission phase, and the
vertical profile showing terrain following against the riverbed.}{fig:overview}

\begin{table}[htbp]\centering\small
\caption{Reference mission results. Navigation error is against simulation
ground truth, which the vehicle never observes.}
\begin{tabular}{lr}
\toprule
\textbf{Metric} & \textbf{Value} \\
\midrule
Mission duration & \SI{\misDuration}{\second} (\misDurationMin{} min) \\
Path length & \SI{\misPath}{\metre} \\
Sonar pings & \misPings{} \\
Soundings mapped & \misSoundings{} \\
Search box coverage & \SI{\misCoverage}{\percent} \\
DVL availability & \SI{\misDvl}{\percent} \\
Navigation error, mean & \SI{\misNavMean}{\metre} \\
Navigation error, maximum & \SI{\misNavMax}{\metre} \\
Navigation drift, fraction of distance travelled & \SI{\misDriftPct}{\percent} \\
Bathymetric map RMSE against truth & \SI{\mapRmse}{\metre} \\
\bottomrule
\end{tabular}
\end{table}

\keyfig{../results/figures/f2_navigation.png}
{Navigation error against the filter's own uncertainty estimate, altitude
tracking, and estimated against true current with both vehicle envelopes
marked.}{fig:nav}

\keyfig{../results/figures/f3_mapping.png}
{Ground truth bathymetry, the map built from simulated sonar during the
mission, and the difference. The residual is near zero over open bed; the
remaining structure is the submerged objects, which stand above the bare bed
and are the subject of Section~\ref{sec:results}.}{fig:mapping}

\subsection{Scour measurement}

\begin{table}[htbp]\centering\small
\caption{Scour measured autonomously against ground truth.}
\begin{tabular}{lccc}
\toprule
\textbf{Pier} & \textbf{Measured depth} & \textbf{True depth} & \textbf{Recovery} \\
\midrule
P2 & \SI{\scourPierTwoDepth}{\metre} & \SI{\scourPierTwoDepthTruth}{\metre} & \SI{\scourPierTwoRec}{\percent} \\
P3 & \SI{\scourPierThreeDepth}{\metre} & \SI{\scourPierThreeDepthTruth}{\metre} & \SI{\scourPierThreeRec}{\percent} \\
\bottomrule
\end{tabular}
\end{table}

Measured scour is systematically shallow by roughly
\SI{\scourMeanAbsErr}{\metre}. The cause is understood: the sonar footprint
averages across the narrow bottom of the scour cone, the map cell is
\SI{0.5}{\metre}, and returns at very shallow grazing angles are rejected as
poorly localised. The bias is consistent and therefore correctable by
calibration, but it is reported here uncorrected.

\subsection{Target detection}

\begin{table}[htbp]\centering\small
\caption{Geometric detection from the bathymetric map. No training data is
used. Small boulders fall below the map cell size and the sonar footprint and
are not separable from bed texture.}
\begin{tabular}{lcc}
\toprule
\textbf{Class} & \textbf{Detected} & \textbf{Present} \\
\midrule
Large vehicle (bus) & \detVehicleLargeHit{} & \detVehicleLargeTot{} \\
Small vehicle (car) & \detVehicleSmallHit{} & \detVehicleSmallTot{} \\
Structural debris & \detStructureHit{} & \detStructureTot{} \\
Boulders & \detBoulderHit{} & \detBoulderTot{} \\
\midrule
\multicolumn{3}{l}{Primary targets: \detPrimary{}/\detPrimaryTotal{},
mean localisation error \SI{\detMeanErr}{\metre}} \\
\bottomrule
\end{tabular}
\end{table}

All vehicles and structural debris present in the scene were found, at a mean
localisation error of \SI{\detMeanErr}{\metre} and a worst case of
\SI{\detMaxErr}{\metre}. For a search and rescue task this is the operationally
meaningful result: it places a diver or recovery crane within a boat length of
the target, in water where neither could otherwise be directed at all.

\subsection{Simulator as a training source: a negative result, stated plainly}

A detector was trained on \synthTrain{} simulated frames with analytically
derived labels and no human annotation, then evaluated on the real ARIS held
out test split.

\begin{table}[htbp]\centering\small
\caption{Detection performance. The first two rows differ only in training
data; the third measures the simulator-trained model inside its own domain.}
\begin{tabular}{llcc}
\toprule
\textbf{Trained on} & \textbf{Tested on} & \textbf{mAP@0.5} & \textbf{mAP@0.5:0.95} \\
\midrule
Real sonar (\num{1627} images) & real test split & \SI{\yoloTestMapFifty}{\percent} & \SI{\yoloTestMap}{\percent} \\
Simulator only (\synthTrain{} images) & real test split & \SI{\synthMapFifty}{\percent} & \SI{\synthMap}{\percent} \\
Simulator only (\synthTrain{} images) & simulated val split & \SI{\synthSelfMapFifty}{\percent} & \SI{\synthSelfMap}{\percent} \\
\bottomrule
\end{tabular}
\end{table}

The third row matters as much as the second. The simulator-trained model
reaches \SI{\synthSelfMapFifty}{\percent} mAP inside its own domain, which
establishes that the simulated imagery is internally consistent, correctly
labelled and genuinely learnable. It is not noise. But that learning does not
carry across to real tank sonar, and the failure is symmetric: run in the
opposite direction, the real-trained detector recovers only
\SI{\simRealRecall}{\percent} of \simRealInstances{} simulated targets across
\simRealScenes{} scenes, against a mean confidence of \simRealConfReal{} on
real test images.

The diagnosis is that the residual gap is dominated by scene content and
capture configuration rather than by the acoustic physics. The public data is
a water tank dense with wall returns, mounting rig and cabling, imaged at a
short range window and a near grazing look, and its classes are centimetre
scale objects whose signature depends on fine structure well below what a
geometric model represents. None of that resembles an open riverbed survey of
metre scale targets.

Two conclusions follow, and both shape the system rather than being excuses.
First, the simulator should not be presented as a drop in replacement for real
training data for fine grained classification. Second, and more usefully, the
capabilities this system actually depends on do not require that transfer at
all. Locating a submerged bus is a geometry problem, solved by the detector of
Section~\ref{sec:results} at \detPrimary{}/\detPrimaryTotal{} with no training
data whatsoever. Where a learned classifier is needed, the relevant question
is not zero-shot transfer but whether simulator pretraining reduces how much
real data must be collected, which is measured next.

\subsection{Data efficiency}

The operational situation is that a limited amount of real sonar can be
collected from a target site, but nothing close to enough to train from
scratch. The useful question is therefore whether pretraining on simulator
output lowers that requirement. Two arms, identical except for
initialisation, were trained on the same $N$ real images and evaluated on the
same untouched real test split: one starting from generic natural-image
weights, one from the simulator-trained weights.

\keyfig{../results/figures/f10_data_efficiency.png}
{Real-data efficiency. Both arms are trained on the same real images and
evaluated on the same real test split, differing only in initialisation.}
{fig:dataeff}

% ======================================================================
\section{Verification and honest limitations}

Every quantitative claim in this report is produced by a script in the
repository and regenerated from the measured result files, so the document
cannot drift out of step with the experiments.

\subsection{What is verified}

\begin{itemize}[leftmargin=*,itemsep=2pt]
\item Ray primitive intersection against analytic values to \num{e-9}.
\item Sonar shadow and target contrast against physical expectation.
\item Vehicle buoyancy, drift under no thrust, and thrust envelope against
closed form drag balance.
\item Navigation error against simulation ground truth over a full mission.
\item Mapped bathymetry, scour and target positions against ground truth.
\item Perception on a real, held out sonar test split.
\end{itemize}

\subsection{What is not}

\begin{itemize}[leftmargin=*,itemsep=2pt]
\item \textbf{No physical vehicle exists.} All vehicle results are simulation.
The hydrodynamic coefficients for the faired hull are engineering estimates and
require CFD and tow tank validation.
\item \textbf{The site is representative, not surveyed.} Dimensions are
consistent with published descriptions of the Mahad collapse and with
IRC:SP:35 practice, but they are modelling assumptions, not measured
bathymetry.
\item \textbf{The sediment attenuation term is an approximation}, not a first
principles result.
\item \textbf{Zero-shot transfer from tank data to open river frames fails.}
This is reported rather than hidden, with the diagnosis above.
\item \textbf{Detection recall is quoted for objects above roughly one metre.}
Smaller debris is below the resolution of the map cell and the sonar
footprint.
\end{itemize}

% ======================================================================
\section{Development pathway}

The present state is a verified simulation and autonomy stack with perception
validated on real sonar data. The path to a field system is:

\begin{table}[htbp]\centering\small
\caption{Development pathway.}
\begin{tabular}{p{0.16\linewidth}p{0.36\linewidth}p{0.40\linewidth}}
\toprule
\textbf{Phase} & \textbf{Work} & \textbf{Evidence produced} \\
\midrule
Now &
Simulation, autonomy, mapping, perception, verification &
This report and the accompanying repository \\
Months 1 to 6 &
CFD and tow tank validation of the faired hull; hardware in the loop with
flight controller and edge compute; indigenous thruster and pressure housing
sourcing &
Measured hydrodynamic coefficients; autonomy running on target hardware \\
Months 7 to 12 &
Integrated vehicle; tank and reservoir trials; DVL and sonar integration &
Wet testing against surveyed targets \\
Months 13 to 24 &
Instrumented river trials with a partner agency; field validation of scour
measurement against diver survey &
Field data at operational current and turbidity \\
\bottomrule
\end{tabular}
\end{table}

The autonomy and perception software is the defensible indigenous content. The
hardware baseline starts from proven components to remove mechanical risk and
is intended to be progressively indigenised across the programme, with the
pressure housing, thrusters and eventually the Doppler velocity log as the
targets in that order.

% ======================================================================
\section*{Reproducing these results}
\addcontentsline{toc}{section}{Reproducing these results}

\begin{verbatim}
python tests/smoke_sonar.py      sonar physics verification
python tests/smoke_control.py    dynamics, estimation, control verification
python eval/thrust_envelope.py   vehicle envelope analysis
python eval/run_mission.py       full autonomous mission
python eval/make_figures.py      report figures
python eval/detect_eval.py       geometric target detection
python eval/sim2real.py          zero-shot transfer test
\end{verbatim}

\end{document}
